\cleardoublepage % memastikan bab baru dimulai di halaman ganjil (kanan)
\chapter{PENDAHULUAN}
\label{chap:pendahuluan}
% --- Latar Belakang ---
\section{Latar Belakang}
Kemajuan internet yang begitu pesat dan cepat menuntut adanya standar dalam tata cara/protokol untuk melaksanakan komunikasi dalam internet agar berbagai jenis perangkat dapat secara mudah terhubung dengan satu sama lain. Salah satu protokol yang cukup terkenal digunakan di internet adalah \textit{Hypertext Transfer Protocol} (HTTP). HTTP merupakan salah satu protokol yang digunakan untuk menghubungkan \textit{client} dan \textit{server} pada \textit{application layer}.

HTTP terus dikembangkan dan ditingkatkan untuk memenuhi kebutuhan jaman yang semakin berubah dengan perubahan yang terbaru adalah HTTP/3.0. HTTP/3.0 telah banyak diimplementasikan sejak HTTP/3.0 masih berupa draf, termasuk oleh perangkat lunak \textit{Web Server}. Akan tetapi, salah satu web server popular, yakni Apache Web Server belum mengimplementasikan HTTP/3.0 secara \textit{native}. Salah satu cara yang sering digunakan oleh penyedia layanan untuk menangani hal tersebut adalah dengan menggunakan \textit{proxy server} yang telah mendukung HTTP/3.0. Akan tetapi, hal ini dapat meningkatkan latensi karena untuk setiap permintaan pengguna harus diproses dua kali. Oleh sebab itu, disini penulis mencoba mengembangkan modul HTTP/3.0 untuk Apache Web Server sehingga penyedia layanan yang menjalankan Apache Web Server dapat mendukung HTTP/3.0 tanpa perlu melalui \textit{proxy server} terlebih dahulu.
% --- Rumusan Masalah ---
\section{Rumusan Masalah}
Hingga saat ini, Apache Web Server belum mendukung HTTP/3.0 secara \textit{native}. Kondisi ini menyebabkan sebagian penyedia layanan yang memanfaatkan \textit{proxy server} yang telah mendukung HTTP/3.0 untuk memberikan dukungan terhadap permintaan pengguna yang menggunakan HTTP/3.0. Akan tetapi hal ini menambah latensi setiap permintaan \textit{client} karena setiap permintaan harus diproses dulu di \textit{proxy server} sebelum nantinya diteruskan ke Apache Web Server. Sebagian penyedia layanan lainnya juga memilih untuk tidak memberi dukungan HTTP/3.0 terlebih dahulu yang menghambat laju adopsi HTTP/3.0 dan kehilangan peluang untuk memanfaatkan fitur-fitur yang ditawarkan HTTP/3.0.

Oleh sebab itu, dibutuhkan suatu modul tambahan untuk Apache Web Server yang mendukung fitur HTTP/3.0. Selain menghilangkan peran \textit{proxy server}, adanya modul ini juga dapat memudahkan penyedia layanan untuk mengadopsi HTTP/3.0.

% --- Tujuan ---
\section{Tujuan}
Hasil yang ingin dicapai dari tugas akhir ini adalah modul HTTP/3.0 untuk Apache Web Server sehingga Apache Web Server dapat menangani permintaan berbasis HTTP/3.0 secara \textit{native}. Untuk memastikan keberhasilan dari tugas akhir ini akan diukur parameter-parameter berikut:
\begin{enumerate}
\item Jumlah permintaan per detik yang dapat diterima \textit{server}
\item Waktu respons rata-rata
\item Waktu respons maksimum (untuk menilai keandalan server saat dibombardir permintaan)
\item Jumlah koneksi aktif yang dapat ditangani \textit{server}
\end{enumerate}

Parameter-parameter ini juga akan dibandingkan dengan pendekatan implementasi HTTP/3.0 via \textit{proxy server} NGINX untuk mengukur efektivitas pengembangan modul HTTP/3.0 \textit{native}.

% --- Batasan Masalah ---
\section{Ruang Lingkup, Asumsi, dan Batasan Masalah}
Dalam pengerjaan tugas ini, modul yang dikembangkan sekurang-sekurangnya harus memenuhi fitur-fitur wajib sesuai dengan ketentuan yang diatur pada \cite{rfc2119}. Selain itu, pengembangan dan pengujian modul ini juga akan dilakukan untuk platform Windows 11 x64. Saat pengujian perbandingan dengan pendekatan yang lain, diasumsikan kondisi perangkat saat pengujian tersebut ekivalen serta diasumsikan juga bahwa tidak ada keunggulan yang dihasilkan karena perbedaan \textit{build tools} antara modul yang dikembangkan dengan NGINX.
\section{Hipotugas akhir}
Tuliskan hipotugas akhir yang akan diuji dalam pelaksanaan tugas akhir. 
Hipotugas akhir ini merupakan pernyataan yang dapat diuji kebenarannya melalui proses pelaksanaan tugas akhir, 
terutama pada tahap evaluasi. Hipotugas akhir yang baik memiliki karakteristik sebagai berikut:
\begin{enumerate}
    \item Spesifik: Hipotugas akhir harus jelas dan spesifik, sehingga dapat diuji dengan cara yang terukur.
    \item Terukur: Hipotugas akhir harus dapat diuji dengan data atau informasi yang dapat diukur, sehingga dapat dinilai kebenarannya.
    \item Relevan: Hipotugas akhir harus relevan dengan topik tugas akhir dan dapat memberikan kontribusi yang signifikan terhadap penyelesaian masalah yang dibahas.
    \item Dapat diuji: Hipotugas akhir harus dapat diuji melalui proses pelaksanaan tugas akhir, terutama pada tahap evaluasi, sehingga dapat dinilai kebenarannya.
\end{enumerate}

Hipotugas akhir diperoleh dari premis-premis yang ada pada latar belakang dan rumusan masalah. 
Oleh karena itu, tuliskan dulu premis-premis yang menentukan hipotugas akhir tersebut sebelum menuliskan hipotugas akhirnya.



% --- Metodologi Pengerjaan TA ---
\section{Metodologi}
Pengerjaan tugas ini mengacu pada dokumen RFC 9114 sehingga sebelum masuk ke fase pengembangan modul akan dikaji terlebih dahulu mengenai fitur-fitur apa yang perlu atau tidak boleh ada dalam HTTP/3.0 sesuai dengan dokumen RFC 9114. Selain itu, proses pengembangan modul akan mengikuti ketentuan pengembangan fitur dalam Apache Web Server, sehingga modul yang dihasilkan dapat digunakan sebagaimana penggunaan modul-modul yang sudah ada.

Pengembangan modul akan dilakukan per fitur secara \textit{incremental}. (gatau mau nulis apa lagi dawg)
% --- Sistematika Penulisan ---
\section{Sistematika Penulisan}
Untuk memudahkan dalam pembacaan proposal, diberikan suatu sistematika penulisan proposal yang berisi penjelasan mengenai konten yang disajikan dalam setiap bab sebagai berikut:
\begin{enumerate}
\item	Bab I Pendahuluan: Berisi latar belakang, rumusan masalah, tujuan, batasan masalah, metodologi, dan sistematika penulisan laporan tugas akhir       .
\item	Bab II Studi Literatur: Berisi kajian pustaka yang relevan dengan topik tugas akhir, termasuk teori-teori, konsep-konsep, dan hasil-hasil penelitian sebelumnya yang                mendukung pelaksanaan tugas akhir.
\item	Bab III Analisis dan Desain: Berisi penjelasan tentang analisis masalah dan kebutuhan pengguna serta berbagai alternatif solusinya.
\item   Bab IV Perancangan: Berisi penjelasan tentang perancangan solusi yang diusulkan, misalnya arsitektur sistem, desain basis data, antarmuka pengguna, dan sebagainya.
\item	Bab V Implementasi: Berisi penjelasan tentang proses implementasi solusi yang diusulkan, termasuk langkah-langkah yang diambil, teknologi yang digunakan, dan hasil yang diperoleh.
\item	Bab VI Evaluasi: Berisi analisis dan evaluasi terhadap solusi yang diimplementasikan, termasuk metode pengujian, hal-hal yang disiapkan sebelum pengujian, hasil pengujian, dan penilaian kinerja.                  
\item   Bab VII Kesimpulan dan Saran: Berisi kesimpulan dari hasil pelaksanaan tugas akhir dan saran untuk pengembangan lebih lanjut.
\item   Daftar Pustaka: Berisi daftar referensi yang digunakan dalam penyusunan laporan tugas akhir.
\item   Lampiran: Berisi dokumen-dokumen pendukung seperti kode program, data lengkap hasil pengujian, rekaman wawancara, dan lain-lain. 
\end{enumerate}